\subsection{THETA3-A8: the completed Taylor tensor is the canonical harmonic jet}\label{A8}

This comparison concerns a fixed polarized abelian variety and its
rigidified theta line bundle.  When the polarized torus and its CM
action in 032 are realized by an archive datum, it identifies the
actual analytic section and its derivatives.  It does not identify
an arbitrary representation in a theta packet with that section.
The distinction between a divided Taylor coefficient and a
differential operator is retained throughout.

\paragraph{The completed theta function is the Appell--Humbert reading.}
First use principal symplectic coordinates.  Write $Z=X+iY$ with
$Y>0$, put $\Lambda=\Z^g+Z\Z^g$, and set
\[
 H(z,w)=z^tY^{-1}\overline w,\qquad
 \widehat\vartheta(Z,z)
       =\exp\left(\frac\pi2z^tY^{-1}z\right)\vartheta(Z,z).
 \tag{A8a}
\]
For $\lambda=a+Zb$ with $a,b\in\Z^g$, the usual theta multiplier
and a direct calculation give
\[
 \widehat\vartheta(Z,z+\lambda)
  =(-1)^{a^tb}
    \exp\left(\pi H(z,\lambda)+\frac\pi2H(\lambda,\lambda)\right)
       \widehat\vartheta(Z,z).
 \tag{A8b}
\]
Indeed, the ordinary multiplier is
$\exp(-\pi i b^tZb-2\pi i b^tz)$.  The new linear term is
$\pi z^t(Y^{-1}\lambda-2ib)=\pi z^tY^{-1}\overline\lambda$.
For the constant term, use
$\lambda-\overline\lambda=2iYb$ to obtain
\[
 \frac\pi2\lambda^tY^{-1}\lambda-\pi i b^tZb
  =\frac\pi2\lambda^tY^{-1}\overline\lambda+\pi i a^tb.
\]
This proves (A8b), including its sign.

For a finite Schwartz input, choose a common lattice level $N$ so
that the positions in its support lie in $N^{-1}\Z^g$ and its
coefficients are invariant under $N\Z^g$.  The same reindexing
argument gives (A8b) on $N\Lambda$.  Equivalently,
$w\mapsto\widehat\vartheta_\varphi(Z,Nw)$ is a section of the
$N^2$-th power of the line bundle on $\C^g/\Lambda$ with the
displayed multiplier.  A nonprincipal polarization is handled in
an integral symplectic lattice with its actual Riemann form; the
same calculation uses its corresponding Hermitian form and
semicharacter.  In the restriction-of-scalars coordinates of 032,
the exponent and Hermitian form are the sums over real embeddings.

The archive's \texttt{AnFr.IsTh}, in
\texttt{Definitions/ThetaValues/ThConst.lean}, is exactly the law
\[
 \theta(z+\lambda)
   =\chi(\lambda)^k
       \exp\left(k\pi H(z,\lambda)
                +\frac{k\pi}{2}H(\lambda,\lambda)\right)\theta(z).
\]
Thus (A8a), with the stated level scaling, is its analytic reading.
Here the polarized lattice, semicharacter, and rigidification must
be matched, and the two readings must represent the same actual
section.  For that section, two nowhere-vanishing trivialization
identifications with the same multiplier have a holomorphic
periodic frame ratio.  This ratio descends to a compact complex
torus, hence is constant, and the rigidification fixes that
constant.  This does not claim that the multiplier determines an
arbitrary section of a higher tensor power.  The $E$-adjointness of
$H$ is the Rosati adjointness of
the polarization and supplies \texttt{AnPt.adj}.  The tangent
frame supplies \texttt{AnPt.der}.

There is a conjugation in the embedding labels that must be kept.
The torus with period $Z$ is naturally
$A=\mathcal F_P^\vee/\iota(\Lambda_{\mathrm{rat}})$, with
$\iota(u)(v)=\psi_P(v,u)$ for $v\in\mathcal F_P$.
In the graph frame this is $\iota(p,q)=Zq-p$, which gives the
displayed period lattice.  Rosati adjointness gives
$\iota(au)(v)=\iota(u)(\overline a v)$.
Consequently the tangent variable dual to the tensor block
$(\mathcal F_P)_\tau$ has the archive tangent label $c\tau$.
The $c\mu$ Taylor-tensor block of 032 therefore corresponds to
the $\mu$ tangent variables in this convention.

\paragraph{What the archive's canonical derivative means.}
The exact definitions are as follows.  In
\texttt{Formalism/LocalToGlobal/ProjectiveLattices/IsThRule.lean},
\texttt{NabV} assigns to a framed point $y$, a section $F$ of
$L^k$, and a monomial $\alpha$ in the $E$-eigenvectors of the
tangent space a scalar $\operatorname{nab}_y^\alpha(F)$.
\texttt{ThGerm.lean}, \texttt{ChartAt} and \texttt{IsNab}, require
that, for $k\geq0$, these scalars be computed from a formal-group
chart, a compatible biextension frame, a normalized primitive,
and a canonical gauge.  They are not arbitrary ordinary
derivatives in an unspecified local generator.

The precise correction is in
\texttt{ExpTiePf.lean}, theorem \texttt{tie\_holds}: if $f$ is the
section's formal coefficient in such a chart, $\phi$ is the formal
exponential with the prescribed tangent frame, and $h_L$ is the
chart's canonical potential, then the corrected exponential germ
is
\[
       g=\bigl(f\exp(kh_L)\bigr)\circ\phi,
 \qquad
       \operatorname{nab}_y^\alpha(F)=\alpha!\,[z^\alpha]g.
 \tag{A8c}
\]
This is the actual sign and factorial convention.  The definition
\texttt{expVal} in \texttt{IsExpChart.lean} includes $\alpha!$.
The analytic chart constructed in
\texttt{ExistsAnalyticBiextensionChartComplexPoint.lean} has
additive formal group and $h=0$, and its germ is the Taylor series
of the Appell--Humbert reading.  The two proof files
\texttt{HarmonicValueCanonicalConnectionEqPolyDeriv.lean} and
\texttt{ThetaValueAnalyticFrameEqPolyDerivFock.lean} use exactly
these data, rather than identifying a raw theta series with an
integral local coefficient.

For completeness, the residual ambiguity has a precise and much
smaller form.  The normalized exponential-frame comparison
\texttt{GaugeUniqSt}, proved in \texttt{GaugeUniqPf.lean}, gives
\[
       g'_b=g_b\exp(-kq),\qquad
 q(z)=\sum_{\tau,i,j}q_{\tau ij}z_{\tau i}z_{c\tau,j},
 \tag{A8d}
\]
where one representative of each conjugate pair is used.
In particular $q$ is homogeneous of degree two; there are no
linear or higher terms.  These restrictions come from the
rigidification and the CM adjointness cocycle in
\texttt{IsExpGauge}, not from a choice to disregard other terms.
The formal exponential is unique for its tangent frame.
For the chart used by \texttt{IsNab}, compatibility under jet
truncation follows from existence of germs for every section of
every nonnegative tensor power: this is
\texttt{GlobalTheory.NablaPinCompare.compat\_of\_germs} in
\texttt{ExistsIsNabNotIsNablaPin.lean}, using properness and
ampleness.  It is therefore legitimate to apply the exponential
and gauge comparison to the actual \texttt{IsNab} data.  All the
archive statements in this paragraph are source inputs read at
the pinned mirror commit; no Lean build or transitive dependency
certificate is claimed here.

\paragraph{A stronger harmonic comparison, with proof.}
Let $z_{\tau i b}$ denote the coordinate in tangent embedding
$\tau$, row $i$, and section column $b$.  Let
\[
 \Delta_{\tau ij}
     =\sum_b
       \frac{\partial^2}
       {\partial X_{\tau i b}\partial X_{c\tau,jb}}.
\]
For any polynomial $P$ satisfying
$\Delta_{\tau ij}P=0$ for every conjugate pair and all rows, and
sections $F_b$ of the same nonnegative power $L^k$, one has
\[
 \sum_\alpha P_\alpha
      \prod_b\operatorname{nab}_y^{\alpha_b}(F_b)
   =P(\partial)\left(\prod_b\theta_b(z_b)\right)\bigg|_{z=0},
 \tag{A8e}
\]
where $\theta_b$ are their Appell--Humbert readings in the matched
tangent frame.  The scalar coefficients can be complex; the
arithmetic version uses the indicated coefficient field.

Here is a direct proof, which also explains why highest-weight
conditions are unnecessary for this particular equality.  Pair a
polynomial and a formal series by
$[P,F]=P(\partial)F(0)$.  On monomials,
\[
       [P,z_aF]=[\partial_{X_a}P,F].
\]
Indeed both sides are zero unless the relevant exponents agree,
and in the remaining case both are the same product of
factorials.  Iteration proves
$[P,Q(z)F]=[Q(\partial_X)P,F]$ for every polynomial $Q$.
Only finitely many terms of an exponential contribute to this
pairing.  Therefore, writing $Q(z)=\sum_b q(z_b)$, (A8d) gives
\[
 [P,\exp(-kQ)F]
      =[\exp(-kQ(\partial_X))P,F]=[P,F],
\]
because $Q(\partial_X)$ is a linear combination of the displayed
$\Delta_{\tau ij}$ and annihilates $P$.  Apply (A8c) in each
chart, use uniqueness of the exponential coordinates, and compare
with the analytic chart.  This proves (A8e).  In particular every
polynomial supported in just one embedding is covered, regardless
of its row weights.

The archive's \texttt{HHa} is more restrictive: by
\texttt{RowDeg.lean}, it is the joint harmonic space together with
lower-row-unipotent invariance and specified degree in every row.
Those extra conditions select highest-weight vectors for the
automorphic construction.  They are not needed for the
coefficient identity just proved.  Equation (A8e) thus extends the
archive's stated harmonic comparison at this precise point; it
does not assert an extension of its automorphic generation or
integrality theorems to arbitrary weights.

\paragraph{The determinant for A6 and all normalization factors.}
If $n$ theta sections live on the same torus and the chosen
embedding has at least $n$ tangent rows, then
\[
       P=\det(X_{\sigma i b})_{1\leq i,b\leq n}
 \tag{A8f}
\]
belongs to \texttt{HHa}.  It is harmonic because it uses no
$c\sigma$ variables.  Its row degrees are one in the first $n$
rows of $\sigma$ and zero elsewhere.  Replacing a selected lower
row by that row plus a multiple of an earlier row preserves its
determinant; operations on an unselected row leave it unchanged.
This proves the required row invariance, including all rows beyond
$n$.  Its coefficients are integers, so it also satisfies the
coefficient condition in \texttt{harmA}.

In the A6 construction the prescribed line types can differ at
compact places.  Their honest common torus is then the product of
the individual polarized tori.  At the active embedding $c\mu$,
the tensor block has dual tangent label $\mu$ by the preceding
calculation.  Write those tangent coordinates as $z_{ij}$, where $i$ is a
cotangent direction and $j$ is the multiplicity-line factor.  The
product theta section is
$\widehat\Theta(z)=\prod_{j=1}^n\widehat\theta_j(z_j)$, and the
required operator is
\[
 \mathcal D=\det(\partial_{z_{ij}}),\qquad
 \mathcal D\widehat\Theta(0)
       =\det\bigl(\partial_{z_{ij}}
                         \widehat\theta_j(0)\bigr)_{i,j}.
 \tag{A8g}
\]
Expand the determinant.  In each summand precisely one derivative
acts on each separate factor, so (A8g) follows from the product
rule with coefficient one.  This polynomial need not satisfy the
row-degree or row-invariance conditions for the full
$n^2$-dimensional embedding block of the product torus.  It does
use only $\mu$ tangent coordinates, so the proved extension (A8e)
applies.  This closes the analytic comparison for the actual
different-type product used in A6, whenever the indicated
polarized product is placed in the archive datum.

To check the factorial separately, write
$\widehat\Theta(z)=\sum_d\langle T_d,z^{\otimes d}\rangle$
as in 032, (4.10).  For homogeneous $P$ of degree $d$, with the
natural symmetric-tensor pairing,
\[
          P(\partial)\widehat\Theta(0)
                 =d!\,\langle P,T_d\rangle.
 \tag{A8h}
\]
For a multiindex $\alpha$ of degree $d$, both sides equal
$\alpha!$ times the coefficient of $z^\alpha$; this proves the
formula.  The double alternation in A6 maps a symmetric product
of $n$ degree-one tensors to the two exterior products with
coefficient one.  Its determinant functional is $n!$ times the
natural symmetric-tensor pairing in (A8h).  Consequently (A8g),
the A6 wedge, and the archive's ordinary differential convention
have precisely the same determinant coefficient.  There is no
extra $n!$ in the wedge itself.

For the archive's translated theta values, its second comparison
theorem retains the additional scalar $M^{-\deg P}$ and uses
\[
 \mathcal F_{q,a}\theta(z)
   =\exp\left(-q\pi H(z,a)-\frac{q\pi}{2}H(a,a)\right)
       \theta(z+a),\qquad q=M^2.
 \tag{A8i}
\]
The translation is $a=M^{-2}\operatorname{mk}(x_2)$ and the
untranslated function is $w\mapsto\theta(Mw)$.  The oddness,
ideal-divisibility and representative conditions on $M,x,x_2$ in
that theorem remain part of its hypothesis.  They are not
consequences of an arbitrary finite Schwartz input.  Likewise,
the final automorphic convention must retain the archive's
complex conjugations and its negative standard additive
character; the analytic equality (A8e) does not silently change
032's positive additive character.

\paragraph{Specialization to the CM elliptic frame in C18.}
For a CM elliptic curve with the invariant coordinate $z$ and
Weierstrass sigma function, the reduced theta function has the
form
\[
 \theta_{\mathrm{red}}(z)
       =\exp(-e_2^*z^2/2)\sigma(z).
 \tag{A8j}
\]
This is the standard reduced-theta normalization used in
Bannai--Kobayashi, Example 1.9, pp.~10--11 of
\url{https://arxiv.org/abs/math/0610163v4}.
If $H(z,a)=h z\overline a$ with $h>0$ and
$\eta(a)=e_2^*a+\pi h\overline a$, direct substitution in (A8i)
gives
\[
 \exp\left(-\pi H(z,a)-\frac\pi2H(a,a)\right)
          \theta_{\mathrm{red}}(z+a)
       =\exp(-e_2^*z^2/2)\,
             e^{-\eta(a)(z+a/2)}\sigma(z+a).
 \tag{A8k}
\]
Thus the common factor for the canonical analytic reading of
C18's tame translated sections is
\[
      q_{\mathrm{can}}(t)
       =\exp\bigl(-e_2^*\lambda(t)^2/2\bigr)
                         \frac{\sigma(\lambda(t))}{t}.
 \tag{A8l}
\]
Equal-degree ratios cancel this entire factor.  For tensor degree
$d$ the exact normalized operator on its integral coefficient is
$q_{\mathrm{can}}^{-d}D(q_{\mathrm{can}}^d\,\cdot)$.
On an elliptic curve with a CM field acting through a single
tangent embedding, the opposite tangent multiplicity is zero.
The quadratic ambiguity (A8d) is therefore zero, and (A8c)
identifies every nonnegative-order canonical jet with the actual
Taylor derivative of the reduced-theta reading.  Harmonicity
does not discard the nonquadratic factor in (A8l): it is the
specified conversion from the integral local generator to the
canonical analytic generator.

\paragraph{Exact archive source locations.}
All these locators refer to the recovered mirror at commit
\texttt{7cbf7a0c32fca1badf05931879a37c3f8a2b03b4}.
The common prefix for \texttt{IsThRule.lean},
\texttt{ThGerm.lean}, \texttt{IsExpChart.lean}, and
\texttt{RowDeg.lean} is
\texttt{Formalism/LocalToGlobal/ProjectiveLattices/}.
The prefix for \texttt{ExpTiePf.lean},
\texttt{GaugeUniq.lean}, and \texttt{GaugeUniqPf.lean} is
\texttt{Formalism/SharpRestriction/Divisibility/GeneratorPowers/}.
The prefix for the three files beginning
\texttt{ExistsAnalyticBiextensionChart},
\texttt{HarmonicValueCanonicalConnection}, and
\texttt{ThetaValueAnalyticFrame}, whose full basenames are given
above, is
\texttt{Development/AbelianVarietyTheory/Proofs/CuspFormEvaluations/}.
The file \texttt{ExistsIsNabNotIsNablaPin.lean} is under
\texttt{Development/AbelianVarietyTheory/Proofs/UnitaryExistence/}.
Finally \texttt{ThConst.lean} is under
\texttt{Development/AbelianVarietyTheory/Definitions/ThetaValues/}.

The arithmetic assertion requires that the chosen imaginary
quadratic CM order, its action, the rigidified polarization and
the invariant differential be the ones in the archive datum.
Neither a general CM field nor an unspecified analytic torus
provides those data automatically.  For the weak tensor-degree
estimate of C18, passage from its raw sigma normalization to
(A8l) additionally requires control of $e_2^*$; equal-degree
ratios do not require it.  If $e_2^*$ is integral, the same Bell
proof extends by allowing the quadratic term: its new weight
inequality is $1/(p-1)\leq2p/(p^2-1)$.  This last observation is
an elementary implication, not an unproved integrality assertion
about $e_2^*$.

\paragraph{The resulting inert-prime specialization.}
Under C19's hypotheses, the required integrality of $e_2^*$ is
proved.  Suppose in addition that the algebraic differential,
rigidification and tame torsion lifts are matched as above, and
that any change of the integral tangent frame has unit determinant.
For a local tensor expression of degree $d$, its canonical
analytic reading is $Q^dv$ with $Q=q_{\mathrm{can}}$ and
$v\in R[[t]]$.  Its order-$r$ canonical local jet is therefore
$D^r(Q^dv)(0)$.  C19 gives the explicit bound
\[
 v_p\bigl(D^r(Q^dv)(0)\bigr)
 \geq\max\left\{0,\frac{r}{p^2-1}
                         -\frac{s_p(r)}{p-1}\right\}.
\]
For equal-degree quotients the common factor cancels and the
bound improves to
\[
 v_p\bigl(D^rv(0)\bigr)
 \geq\max\left\{0,
          \left\lfloor\frac{rp}{p^2-1}\right\rfloor-1\right\}.
\]
For actual archive sections of nonnegative power these are the
jets in (A8c).  For local ratios and negative tensor degree, they
are the induced local connection expressions, obtained by the
product and quotient rules.  They make no assertion about the
unconstrained negative-power values of an arbitrary
\texttt{NabV} satisfying only \texttt{IsNab}.

The level normalization can also be checked exactly.  For the
degree-one reduced theta function and integer $M$, linearity of
$H$ gives
\[
 \mathcal F_{M^2,a}\bigl(z\mapsto\theta_{\mathrm{red}}(Mz)\bigr)(z)
      =\bigl(\mathcal F_{1,Ma}\theta_{\mathrm{red}}\bigr)(Mz).
\]
An order-$r$ derivative on the right contributes $M^r$, which
cancels the archive's $M^{-r}$ in (A8i).  Thus, when $M$ and the
relevant torsion order are prime to $p$, the displayed C19 bounds
apply to this normalized value without an additional level
denominator.  This is an elliptic specialization with specified
normalizations, not an assertion that the higher-dimensional
032 torus has such an elliptic realization.

\framework{The comparison fixes a rigidified line bundle, its
formal completion, and the correspondence between an algebraic
section and its analytic reading.  It is compatible with the
correspondence and polarization viewpoint in Grothendieck's
\emph{cote44}, scans 7--8, and with the dimension and Lefschetz
discussion in \emph{cote16}, scans 43--46.  No algebraic inverse
Lefschetz operator or Chow projector is inferred from equality of
analytic jets.}
